\documentclass[tikz,border=4pt]{standalone}
% Preámbulo común de los diagramas del Objetivo 3 (paleta de main.tex).
\usepackage[T1]{fontenc}
\usepackage[utf8]{inputenc}
\usepackage{lmodern}
\usepackage[spanish]{babel}
\usetikzlibrary{babel,arrows.meta,positioning,fit,backgrounds,shapes.multipart,calc}
\definecolor{dPersonaF}{HTML}{FFE6CC}\definecolor{dPersonaB}{HTML}{D79B00}
\definecolor{dSistemaF}{HTML}{DAE8FC}\definecolor{dSistemaB}{HTML}{6C8EBF}
\definecolor{dAgenteF}{HTML}{D5E8D4}\definecolor{dAgenteB}{HTML}{82B366}
\definecolor{dDatosF}{HTML}{E1D5E7}\definecolor{dDatosB}{HTML}{9673A6}
\definecolor{dClinicoF}{HTML}{FFF2CC}\definecolor{dClinicoB}{HTML}{D6B656}
\definecolor{dExternoF}{HTML}{F5F5F5}\definecolor{dExternoB}{HTML}{666666}
% Entidad: encabezado coloreado y columnas clave. PK subrayada, FK en cursiva.
\tikzset{
  ent/.style 2 args={rectangle split, rectangle split parts=2, draw=#2, line width=0.7pt,
    rectangle split part fill={#1,white}, rounded corners=1.5pt, font=\footnotesize,
    inner sep=3pt, minimum width=2.9cm},
  ext/.style={draw=dExternoB, dashed, fill=dExternoF, rounded corners=1.5pt, font=\footnotesize\ttfamily,
    inner sep=4pt, minimum width=2.4cm},
  rel/.style={draw=black!65, line width=0.6pt},
  card/.style={font=\scriptsize, inner sep=1.5pt, fill=white},
}
\newcommand{\pk}[1]{\underline{#1}}
\newcommand{\fk}[1]{\textit{#1}}
\newcommand{\cols}[1]{\begin{tabular}{@{}l@{}}#1\end{tabular}}
\newcommand{\tname}[1]{\textbf{\ttfamily #1}}
% Relación 1:N con cardinalidades junto a cada extremo.
\newcommand{\unoN}[3][]{\draw[rel,#1] (#2) -- node[card, pos=0.12]{1} node[card, pos=0.88]{N} (#3);}
\newcommand{\unoUno}[3][]{\draw[rel,#1] (#2) -- node[card, pos=0.12]{1} node[card, pos=0.88]{1} (#3);}
\newcommand{\leyenda}{\node[font=\scriptsize, text=black!70, anchor=north west] at (current bounding box.south west)
  {\underline{subrayado}: clave primaria \quad \textit{cursiva}: clave foránea \quad recuadro gris: tabla de otro grupo};}

\begin{document}
\begin{tikzpicture}[
  font=\footnotesize, every node/.append style={execute at begin node={\hyphenpenalty=10000\exhyphenpenalty=10000}},
  paso/.style={draw, rounded corners=2pt, line width=0.7pt, align=center, text width=#1, minimum height=0.8cm, inner sep=3pt},
  paso/.default=3.6cm,
  nota/.style={draw=black!45, dashed, rounded corners=2pt, align=center, font=\scriptsize, inner sep=3pt, fill=white},
  carril/.style={draw=black!25, fill=black!3, rounded corners=3pt},
  titulo/.style={font=\scriptsize\bfseries, anchor=south, text=black!70},
  f/.style={-{Stealth[length=2mm]}, line width=0.65pt, draw=black!70},
  e/.style={font=\scriptsize\itshape, fill=white, inner sep=1pt}
]
% Carril 1: WhatsApp (Meta)
\node[paso=3.2cm, fill=dPersonaF, draw=dPersonaB] (m1) at (0,0) {Paciente envía un mensaje};
\node[paso=3.2cm, fill=dExternoF, draw=dExternoB] (m2) at (0,-1.3) {Meta envía el evento al webhook};
\node[nota, text width=3.0cm] (m3) at (0,-3.9) {Sin respuesta 200 a tiempo, Meta reintenta el envío};
\node[paso=3.2cm, fill=dPersonaF, draw=dPersonaB] (m4) at (0,-8.9) {El paciente recibe la respuesta};
% Carril 2: recepción HTTP
\node[paso, fill=dAgenteF, draw=dAgenteB] (h1) at (4.4,-1.3) {Verifica la firma HMAC SHA-256 con el secreto de la aplicación};
\node[paso, fill=dAgenteF, draw=dAgenteB] (h2) at (4.4,-2.6) {Programa la tarea y responde 200 de inmediato};
\node[nota, text width=2.4cm] (h0) at (4.4,0) {Firma inválida: 401};
\node[nota, text width=3.4cm] (h3) at (4.4,-7.6) {Error: descarta la conversación y envía el aviso de problema técnico con el contacto de recepción};
% Carril 3: tarea en segundo plano
\node[paso=4.2cm, fill=dAgenteF, draw=dAgenteB] (b1) at (9.3,-2.6) {Extrae teléfono, identificador y tipo (texto, botón, lista u otro)};
\node[paso=4.2cm, fill=dAgenteF, draw=dAgenteB] (b2) at (9.3,-3.9) {Descarta el identificador ya procesado (últimos 1000)};
\node[paso=4.2cm, fill=dAgenteF, draw=dAgenteB] (b3) at (9.3,-5.1) {Candado por teléfono: un mensaje a la vez};
\node[paso=4.2cm, fill=dAgenteF, draw=dAgenteB] (b4) at (9.3,-6.3) {Recupera la conversación; expira a los 30 min};
\node[paso=4.2cm, fill=dSistemaF, draw=dSistemaB, line width=1pt] (b5) at (9.3,-7.6) {\textbf{Enrutador}: máquina de estados, modelo y API interna};
\node[paso=4.2cm, fill=dAgenteF, draw=dAgenteB] (b6) at (9.3,-8.9) {Envía las respuestas por la Graph API};
\node[paso=4.2cm, fill=dAgenteF, draw=dAgenteB] (b7) at (9.3,-10.1) {Guarda la conversación y libera el candado};
% Carriles
\begin{scope}[on background layer]
  \draw[carril] (-1.95,0.6) rectangle (1.95,-10.75);
  \draw[carril] (2.2,0.6) rectangle (6.6,-10.75);
  \draw[carril] (6.85,0.6) rectangle (11.75,-10.75);
\end{scope}
\node[titulo] at (0,0.6) {WhatsApp (Meta)};
\node[titulo] at (4.4,0.6) {Recepción del webhook};
\node[titulo] at (9.3,0.6) {Tarea en segundo plano};
% Flujo
\draw[f] (m1) -- (m2);
\draw[f] (m2) -- (h1);
\draw[f] (h1) -- (h2);
\draw[f, dashed] (h1) -- (h0);
\draw[f] (h2) -- node[e, above]{tarea} (b1);
\draw[f, dashed] (h2.west) -| node[e, pos=0.25, above]{200} (m3.north);
\foreach \a/\b in {b1/b2, b2/b3, b3/b4, b4/b5, b5/b6, b6/b7} \draw[f] (\a) -- (\b);
\draw[f] (b6) -- (m4);
\draw[f, dashed] (b5) -- (h3);
\end{tikzpicture}
\end{document}
